\section{Discusión de resultados}
\label{sec:discusion}

La localización depende de la estructura del sistema y de la respuesta que
buscamos. En los sistemas de Lur'e, el balance armónico relaciona la
frecuencia y la amplitud con la transferencia del bloque lineal. Cuando la
ganancia requerida pertenece al rango de la función descriptiva, esa
relación proporciona una semilla. La continuación permite transportarla
hasta el campo de interés. En \(q=1\), esta construcción recupera el método
de Leonov--Kuznetsov \cite{LeonovKuznetsov2013,KuznetsovEtAl2017}.

Las obstrucciones de Kalman--Fitts y PLL muestran por qué necesitamos rutas
complementarias. En Kalman--Fitts, la ganancia que exige el cierre directo
es negativa y queda fuera del rango de la función descriptiva de \(\tanh\).
El mapa de conmutación y la homotopía signo--tanh permiten localizar el
ciclo por otra vía. En PLL, la parte imaginaria de la transferencia
conserva su signo a frecuencia positiva; la continuación de la ganancia
de lazo y el retorno sobre el cilindro aprovechan la geometría de fase
que el cierre directo no resuelve.

En PLL localizamos numéricamente un ciclo de cuenca oculta: las 96
condiciones próximas a los equilibrios no alcanzan el ciclo y sus destinos
coinciden en la comparación Python--Julia. MAVPD presenta un ciclo periódico
externo para \(\xi=3.1\), con \(0/108\) contactos locales, y un régimen
caótico para \(\xi=2.85\), \(\gamma=0.15380379839949113\), con
\(0/112\) contactos. La comparación independiente coincide en las 112
decisiones del segundo caso. El retorno periódico del primer caso permite
distinguirlo de la cuasiperiodicidad atribuida en la referencia. Al cambiar
a \(\xi=3.5\), aparecen 24 contactos: la respuesta pasa a ser accesible
desde las vecindades ensayadas. Este control muestra que la clasificación
depende de los parámetros y de la cuenca, no sólo de la forma visual de
la trayectoria.

Los puntos perpetuos también distinguen destinos dentro de un mismo
modelo. En las ocho repeticiones de MAVPD, la ruta PP produce 91 contactos
interiores de 96 semillas y cinco transitorios; la función descriptiva
reparte sus contactos entre las referencias interior y exterior. Una
comparación restringida al atractor externo perdería esa diferencia.
En PLL, PP obtiene seis contactos por cada doce semillas en todas las
repeticiones, mientras que el muestreo uniforme varía entre cuatro y siete.
La estabilidad de una ruta y su cobertura son propiedades diferentes; los
resultados no establecen una ventaja uniforme de un generador.

Para \(0<q<1\), la construcción frecuencial y la evolución causal cumplen
funciones distintas. El operador de Liouville--Weyl proporciona un predictor
armónico estacionario; el PVI de Caputo conserva la contribución de la
historia desde un terminal inferior fijo. En Chua PWL, la semilla sesgada
se transporta a través de 21 etapas con una sola historia. La respuesta
autónoma tiene una frecuencia dominante consistente con la reconstrucción
de los parámetros de Danca \cite{Danca2017}. El reinicio independiente
reproduce esa frecuencia, con una diferencia máxima de entropía espectral
de \(2.2243710\times10^{-4}\). La recurrencia ADM local de Wu
\cite{Wu2023ArctanChua}, en cambio, reinicia la memoria en cada intervalo
y produce dos respuestas simétricas regulares. La coincidencia de una
figura o de una frecuencia no identifica estos dos problemas de evolución.

\ifHAFOBasinMaterial
En el sondeo PWL, las \(7200\) condiciones del dominio local
\(r\leq0.01\) no alcanzan la referencia. Los radios \(0.03\) y \(0.1\)
añaden \(6600\) condiciones sin contacto y el radio \(0.3\) produce
37 contactos de \(3600\). Así observamos un acceso a la respuesta al
alejarnos de los equilibrios. En c590, los conteos son \(0/8400\) hasta
\(r=0.3\), \(22/2400\) en \(r=1\) y \(588/2700\) en \(r=2\),
con 131 divergencias en este último bloque. Estos sondeos describen el
acceso muestreado; no fijan una frontera exacta entre radios que no se
han explorado de forma continua.
\fi

En c590, la sensibilidad a la discretización afecta la interpretación de
esos contactos. El contraste L1--PECE de memoria completa sólo coincide en
siete de 18 comparaciones de destino, y cuatro de seis PVI cambian alguna
etiqueta al variar método o paso. La semilla de referencia con PECE y
\(h=0.0025\) alcanza la norma de parada en \(t=176.05\). Por ello,
el comportamiento aperiódico de una trayectoria a paso fijo no permite
concluir una cuenca robusta respecto del refinamiento. El caso aporta
información sobre sensibilidad numérica y sobre la necesidad de contrastar
la trayectoria antes de interpretar sus mapas como cuencas asintóticas.

\section{Conclusiones}
\label{sec:conclusiones}

La representación de Lur'e, el balance armónico y la continuación forman
una ruta de localización que se complementa con la geometría del campo.
Las derivaciones establecen cuándo una ganancia y una semilla son
compatibles con las ecuaciones; las integraciones y los sondeos establecen
qué respuesta alcanzan bajo las condiciones ensayadas. Esta separación
permite interpretar tanto los éxitos como las obstrucciones sin exigir
que un mismo generador funcione en todas las familias.

Los resultados enteros de PLL y MAVPD proporcionan ejemplos numéricos de
atractores ocultos periódicos y caóticos, con sondeos próximos a todos los
equilibrios y comparaciones independientes de destino. Kalman--Fitts amplía
la localización a un caso en el que falla el cierre armónico directo.
Los controles MAVPD \(\xi=3.5\) y Lorenz generalizado muestran, a su vez,
que una semilla algebraicamente válida puede alcanzar una respuesta
autoexcitada o aproximarse a un equilibrio.

La adaptación a Caputo exige conservar la historia durante la continuación.
El reinicio puntual cambia el PVI, incluso en una ecuación lineal escalar.
La ruta PWL sesgada produce una oscilación regular y un dominio muestreado
sin contactos locales; el contraste c590 muestra que el paso y el método
pueden cambiar el destino calculado. El análisis de memoria y de error
forma, por tanto, parte de la metodología de localización.

Quedan problemas matemáticos concretos: justificar la persistencia de las
respuestas fraccionarias bajo refinamiento, caracterizar su atracción en
un espacio de historias y construir regiones invariantes que expliquen la
separación de cuencas. El ciclo de
\cref{subsec:guide-hidden-cycle-topology} muestra cómo un argumento
topológico puede completar la localización con una separación demostrada.
Extender esa relación a los sistemas con memoria requiere controlar la
historia y la aproximación del operador, además de sus proyecciones
finitas en el espacio de estados.
